%------------------------------------------------------------------------------%
\begin{question}
  Determine a posição relativa entre as retas
  \[
    r\colon
    \begin{cases}
      x =  1 +  t \\
      y =  2 + 2t \\
      z = -1 + 3t
    \end{cases}
  \qquad\qquad
    s\colon
    \begin{cases}
      x & = 4 + 2\tau \\
      y & = 1 + 4\tau \\
      z & = 2 + 6\tau
    \end{cases}
  \]
\end{question}

%------------------------------------------------------------------------------%
\begin{resolution}{Solução}
\begin{columns}[t]
\column{0.3\textwidth}
  \onslide<+->{Os vetores diretores são}
  \[
    \onslide<+->{\vec{v} = (1, 2, 3)}
  \]
  \[
    \onslide<+->{\vec{w} = (2, 4, 6)}
  \]
  \onslide<+->{Como
  \[
    \vec{w} = 2\vec{v}
  \]}
  \onslide<+->{as retas são paralelas}
\column{0.3\textwidth}
  \onslide<+->{Verificar se
  \[
    P = (1, 2, -1) \in s
  \]}
  \vspace{-\baselineskip}
  \begin{align*}
    \onslide<+->{x     & = 4 + 2\tau    \\}
    \onslide<+->{1     & = 4 + 2\tau    \\}
    \onslide<+->{2\tau & = -3           \\}
    \onslide<+->{\tau  & = \frac{-3}{2}}
  \end{align*}
\column{0.3\textwidth}
  \begin{align*}
    \onslide<+->{y & = 1 + 4\tau         \\}
    \onslide<+->{  & = 1 + 4\frac{-3}{2} \\}
    \onslide<+->{  & = -5 \\}
    \onslide<+->{  & \neq 2}
  \end{align*}
  \onslide<+->{portanto, \(P\notin s\)}

  \medskip
  \onslide<+->{As retas são\\ paralelas distintas}
\end{columns}
\end{resolution}

%------------------------------------------------------------------------------%
\endinput
